% Source PDF: source/普通高考/2009/2009浙江理.pdf, page 1 (question 6).
% Node-edge list: 开始 -> k=0 -> S=0 -> S<100?;
% 是 -> S=S+2^S -> k=k+1 -> loop to the decision; 否 -> 输出 k -> 结束.
% The loop is a single orthogonal return path with one arrow entering the decision stem.
\begin{tikzpicture}[
    >=Stealth,
    x=1cm,
    y=0.86cm,
    line width=0.45pt,
    line cap=round,
    line join=round,
    flowtext/.style={anchor=center,align=center,
        execute at begin node={\setlength{\parindent}{0pt}}},
    every node/.style={flowtext,font=\normalsize},
    flowbox/.style={draw,flowtext,minimum height=0.68cm,inner xsep=4pt,inner ysep=2pt},
    terminal/.style={flowbox,rounded corners=0.18cm,minimum width=1.30cm},
    process/.style={flowbox,minimum width=1.85cm},
    decision/.style={draw,diamond,flowtext,aspect=2.9,minimum width=3.90cm,
        minimum height=1.05cm,inner sep=0pt},
    io/.style={draw,trapezium,flowtext,trapezium left angle=72,
        trapezium right angle=108,minimum width=1.95cm,minimum height=0.76cm,
        inner xsep=4pt,inner ysep=2pt},
    branchlabel/.style={flowtext,inner sep=0pt},
]
    \path[use as bounding box] (-3.55,-0.55) rectangle (4.70,10.55);

    \node[terminal] (start) at (0,10.05) {开始};
    \node[process] (setk) at (0,8.72) {$k=0$};
    \node[process] (sets) at (0,7.42) {$S=0$};
    \node[decision] (check) at (0,6.00) {$S<100?$};
    \node[process,minimum width=2.80cm] (update) at (0,4.35) {$S=S+2^S$};
    \node[process,minimum width=2.35cm] (inc) at (0,2.82) {$k=k+1$};
    \node[io,minimum width=1.95cm] (output) at (3.45,4.55) {输出 $k$};
    \node[terminal] (finish) at (3.45,2.95) {结束};

    \draw[->] (start.south) -- (setk.north);
    \draw[->] (setk.south) -- (sets.north);
    % Main stem and return edge share one merge point above the decision.
    \coordinate (merge) at (0,6.90);
    \draw[-] (sets.south) -- (merge);
    \draw[->] (merge) -- (check.north);
    \draw[->] (check.south) -- node[branchlabel,left=3pt] {是} (update.north);
    \draw[->] (update.south) -- (inc.north);
    % Keep a visible downward shaft before the arrowhead enters 输出 k.
    \draw[-] (check.east) -- node[branchlabel,above=2pt] {否} (3.45,6.00)
        -- (3.45,5.35);
    \draw[->] (3.45,5.35) -- (output.north);
    \draw[->] (output.south) -- (finish.north);

    % Return path from k=k+1 enters the main stem above the decision.
    \draw (inc.south) -- (0,1.85) -- (-2.75,1.85) -- (-2.75,6.90);
    \draw[->] (-2.75,6.90) -- (merge);
\end{tikzpicture}
